\setchoices{label=(\Alph*)}

\begin{question}[index=1]
设 $g(x)=2^{\frac1{x-1}}$，$f(x)=\begin{cases}\dfrac{(x-1)[1-g(x)]}{1+g(x)},&x\ne1,\\0,&x=1,\end{cases}$ 则
\begin{choices}[columns=2]
\item $f(x)$ 在 $x=1$ 处不连续
\item $f(x)$ 在 $x=1$ 处可导
\item $f(x)$ 在 $x=1$ 处取极大值
\item $f(x)$ 在 $x=1$ 处取极小值
\end{choices}
\end{question}
\QuestionBreak

\begin{question}[index=2]
设函数 $y=y(x)$ 由方程 $x^2+y^2-xy=1$ 确定，曲线 $L$ 为 $y=y(x)$ 介于 $(-1,0)$ 和 $(0,1)$，且位于第二象限的一段，则 $\displaystyle\int_L\frac{(y^2-xy)(2y-x)}{\sqrt{5-3xy}}\dd s=$
\begin{choices}[columns=4]
\item $\dfrac12$
\item $\dfrac13$
\item $\dfrac23$
\item $1$
\end{choices}
\end{question}
\QuestionBreak

\begin{question}[index=3]
$\displaystyle\int_0^1\dd x\int_0^1\frac{1}{1+\max\{x^2,y^2\}}\dd y=$
\begin{choices}[columns=4]
\item $\eu^2$
\item $\ln2$
\item $\eu^3$
\item $\ln3$
\end{choices}
\end{question}
\QuestionBreak

\begin{question}[index=4]
设有定积分 $\displaystyle I_1=\int_0^{2\pi}\eu^{-x^2}\sin x\dd x$，$\displaystyle I_2=\int_0^{2\pi}\eu^{-x^2}\cos x\dd x$，则有
\begin{choices}[columns=4]
\item $I_1>0,I_2>0$
\item $I_1>0,I_2<0$
\item $I_1<0,I_2>0$
\item $I_1<0,I_2<0$
\end{choices}
\end{question}
\QuestionBreak

\begin{question}[index=5]
设 $3$ 阶方阵 $A=(\alpha_1,\alpha_2,\alpha_3)$，$\alpha_1,\alpha_2,\alpha_3$ 为三维列向量，$B=(\alpha_1,\alpha_2+2\alpha_3,\alpha_3)$，$B^*$ 为 $B$ 的伴随矩阵，若 $|A|=2$，则 $B^*A=$
\begin{choices}[columns=4]
\item $\mat{2&0&0\\0&2&0\\0&4&2}$
\item $\mat{2&0&0\\0&2&0\\0&-4&2}$
\item $\mat{2&0&0\\0&2&4\\0&0&2}$
\item $\mat{2&0&0\\0&2&-4\\0&0&2}$
\end{choices}
\end{question}
\QuestionBreak

\begin{question}[index=6]
设 $\alpha_1=(1,1,2)^{\mathrm T}$，$\alpha_2=(1,0,-1)^{\mathrm T}$，$\alpha_3=(2,0,a)^{\mathrm T}$，$\alpha_4=(0,1,b)^{\mathrm T}$，若 $\alpha_4$ 可由 $\alpha_1,\alpha_2,\alpha_3$ \allowbreak\mbox{线性表示}，但 $\alpha_3$ 不能由 $\alpha_1,\alpha_2,\alpha_4$ 表示，则 $a,b$ 的取值为
\begin{choices}[columns=4]
\item $a\ne-2,b\ne3$
\item $a=-2,b\ne3$
\item $a\ne-2,b=3$
\item $a=-2,b=3$
\end{choices}
\end{question}
\QuestionBreak

\begin{question}[index=7]
设 $3$ 阶方阵 $A$ 为正交实对称矩阵，且 $|A|<0$，$\alpha_1=(0,\sqrt2+1,-1)^{\mathrm T}$ 是 $A$ 的属于\allowbreak\mbox{单特征值 $-1$} 的特征向量，$A^*$ 为 $A$ 的伴随矩阵，若方阵 $B$ 满足 $A^*B=|A|E+A^*$，则 $Bx=0$ 的一个\allowbreak\mbox{基础解系}为
\begin{choices}[columns=2]
\item $\mat{0\\1\\\sqrt2+1}$
\item $\mat{0\\\sqrt2+1\\-1}$
\item $\mat{0\\1\\\sqrt2+1},\mat{0\\\sqrt2+1\\-1}$
\item $\mat{1\\0\\0},\mat{0\\1\\\sqrt2+1}$
\end{choices}
\end{question}
\QuestionBreak

\begin{question}[index=8]
设随机变量 $X,Y$ 满足 $X\ge Y$，分布函数分别为 $F_X(x),F_Y(x)$，联合分布函数为 $F(x,y)$，则
\begin{choices}[columns=2]
\item $F_X(x)\ge F_Y(x),F(x,x)\le F_X(x)$
\item $F_X(x)\le F_Y(x),F(x,x)<F_X(x)$
\item $F_X(x)\ge F_Y(x),F(x,x)\le F_Y(x)$
\item $F_X(x)\le F_Y(x),F(x,x)\le F_Y(x)$
\end{choices}
\end{question}
\QuestionBreak

\begin{question}[index=9]
在长度为 $2$ 的细绳上均匀随机选取一点截开，从一端起的第一段围成圆，另一段围成正方形。记 $X$ 表示圆的直径，$Y$ 表示正方形的面积，则 $\operatorname{Cov}(X,Y)=$
\begin{choices}[columns=4]
\item $-\dfrac1{12\pi}$
\item $-\dfrac1{96\pi}$
\item $-\dfrac1{24\pi}$
\item $-\dfrac1{48\pi}$
\end{choices}
\end{question}
\QuestionBreak

\begin{question}[index=10]
设变量 $X_1,X_2,X_3$ 为来自正态总体 $X\dist N(0,3)$ 的简单随机样本，则下列结论正确的个数是
\par\circlednumber{1}\quad $\dfrac13(X_1^2+X_2^2+X_3^2)\dist\chi^2(3)$
\par\circlednumber{2}\quad $\dfrac19(X_1+X_2+X_3)^2\dist\chi^2(1)$
\par\circlednumber{3}\quad $\dfrac{X_1+X_2}{|X_1-X_2|}\dist t(1)$
\par\circlednumber{4}\quad $\dfrac{2X_1^2}{X_2^2+X_3^2}\dist F(1,2)$
\begin{choices}[columns=4]
\item $1$
\item $2$
\item $3$
\item $4$
\end{choices}
\end{question}
